\section{问题重述}

\subsection{问题背景}

无线电干扰源的定位与清除是无线电频谱管理的一项核心任务，其对象是由非法用频、设备杂散辐射与工业设施电磁泄漏等原因产生的干扰源，适用于通信、广播、航空导航与物联网等电磁环境复杂的场景。

该任务按初步监测、区域性排查与精准定位三个阶段开展。初步监测与区域性排查依托固定监测站、移动监测车与无人机完成大范围信号采集，确定干扰源的目标分布区域；精准定位阶段由技术人员携带便携式无线电方位监测测向机进入目标区域徒步排查，测向机依据场强变化推断干扰源的方位，交会定位法由多个检测点的示向度确定干扰源的位置。示向度读数受当地电磁环境与测向精度所限而带有误差，干扰源的最大可被有效接收半径处于 1000 米至 1500 米之间，两者共同决定定位精度与作业效率。

在辐射、有毒有害等人员进入受限的危险现场，搭载测向机、光学探测仪与激光枪的机器狗代替技术人员开展作业，但机器狗作业面临信息与时间双重约束。由多次读数确定干扰源的位置、由有限次检测确定检测点的选取，并据此规划覆盖目标区域的搜索与清除顺序使总时间最短，构成系统的核心建模问题，其求解对提高无线电频谱管理效率、保障危险环境作业安全具有重要意义，为机器狗自动定位清除系统的研制与优化提供依据。

\subsection{问题要求}

本研究结合题设的干扰源分布条件与设备参数，建立数学模型研究以下四个问题：

问题 1 给定若干检测点坐标与某干扰源关于这些检测点的示向度，给出交会定位所得多边形定位区域的直径算法，并判定以该直径为直径的圆对该区域的覆盖关系。

问题 2 已知某全向干扰源在一个检测点处测得的示向度，给出第二检测点的选择策略与候选区域，使该干扰源的定位精度达到最优。

问题 3 目标区域内的干扰源全部为全向干扰源，总数处于 10 至 16 个之间且具体个数未知，机器狗自坐标原点出发，测向机初始频道为 1。制定自动搜索定位与清除的策略及算法，在清除全部干扰源的前提下使定位清除总时间最短，并通过模拟器演练测试给出被清除干扰源个数与平均定位清除时间的统计结果。

问题 4 目标区域内的干扰源包含全向干扰源与定向干扰源，定向干扰源的个数与方向均未知，其余条件与问题 3 相同。给出该条件下的检测与清除策略及算法，并通过模拟器演练测试检验其有效性。
